\begin{center}
    \begin{large}
        \textbf{Abstract}
    \end{large}
    \begin{justify}
        \begin{normalsize}
            Glass fibre-reinforced polymer composites are widely used in the
             wind turbine manufacturing industry, and are commonly manufactured 
             using general-purpose unsaturated polyester resin thermosets.
            However, these resins are obtained from oil-derived aromatic
             monomers making them both unsustainably sourced and difficult
             to recycle.
            Bio-based and biodegradable alternatives exist but have not yet
             seen widespread commercial success due to their relatively
             high cost and poor mechanical properties. 
            This paper seeks to address the question of whether bio-based 
             branching thermoset polyesters can be optimised using Bayesian 
             optimisation such that properties correlated with mechanical
             performance, including extent of polymerisation, 
             might be improved upon.
            This is achieved by selecting a bio-based and biodegradable
             polymer (poly(glycerol citrate itaconate)), identifying a
             desirable objective variable (mass loss under esterification
             as proxy for extent of polymerisation)
             and some tunable predictor variables (monomeric stoichiometry),
             and then applying sequential Bayesian techniques to
             uncover well-performing samples in an experimentally efficient
             manner.
            A Bayesian technique employing the probability of improvement
             acquisition function discovered the experimental global optimum
             stoichiometry producing a mass loss of 20.77\% in curing 
             polyester samples- this represents 72\% of the theoretical 
             maximum, indicating a high extent of polymerisation.
            Bayesian techniques employing both expected improvement and
             probability of improvement acquisition functions were simulated
             10,000 times across five different emulations of the practical
             experiment's parameter space, respectively exhibiting 
             probabilities of discovering the hypothetical `region of global
             optimality' of 65.0\% and 23.8\%.
            We conclude that Bayesian optimisation techniques represent a 
             practically deployable and sample-efficient means by which 
             desirable properties of bio-based thermosets may be improved.
        \end{normalsize}
    \end{justify}
\end{center}
